% !PRINTED 173
\weberpage{173}

\[
\mu_1 + \mu_2 + \cdots + \mu_m = \nu_1 + 2\nu_2 + \cdots + (n-1)\nu_{n-1},
\]
und danach kann der Ausdruck für $\sigma_m$ so dargestellt werden:
\weq{15}{
\sigma_m = \overset{\mu}{\textstyle\sum} s_{\mu_1+\mu_2+\ldots+\mu_m}\, \alpha_{\mu_1}\alpha_{\mu_2}\cdots\alpha_{\mu_m}.
}
Beispielsweise wird
\[
\sigma_1 = \alpha_0 s_0 + \alpha_1 s_1 + \cdots + \alpha_{n-1} s_{n-1}
\]
\weq{16}{
\begin{aligned}
\sigma_2 &= \overset{i,k}{\textstyle\sum} s_{i+k}\, \alpha_i \alpha_k \\
\sigma_3 &= \overset{h,i,k}{\textstyle\sum} s_{h+i+k}\, \alpha_h \alpha_i \alpha_k,
\end{aligned}
}
worin $h$, $i$, $k$ die Reihe der Zahlen $0, 1, 2 \ldots n-1$ durchlaufen.

\begin{center}
\S.\,53.

\gesperrt{Anwendung auf die cubischen und biquadratischen Gleichungen.}
\end{center}

Das Ziel, das schon \gesperrt{Tschirnhausen} bei dieser Transformation im Auge gehabt hat, bestand darin, durch Bestimmung der Substitutionsco\"efficienten $\alpha_0, \alpha_1 \ldots \alpha_{n-1}$ in der umgeformten Gleichung $\Phi(y) = 0$ so viele Glieder zum Verschwinden zu bringen, dass die Gleichung lösbar wird. Es gelingt dies leicht bei den Gleichungen dritten und vierten Grades, wenn es auch nicht einfach ist, und nicht ohne weitläufige Rechnungen möglich scheint, die sonst bekannten Formeln auf diesem Wege abzuleiten.

Nehmen wir zunächst die cubische Gleichung
\weq{1}{
x^3 + a x^2 + b x + c = 0,
}
die wir durch die Substitution
\weq{2}{
y = \alpha + \beta x + \gamma x^2
}
umformen. Um die Gleichung für $y$
\weq{3}{
y^3 + A y^2 + B y + C = 0
}
auf eine reine cubische Gleichung zurückzuführen, haben wir die Verhältnisse $\alpha:\beta:\gamma$ aus den beiden Gleichungen
\weq{4}{
A = 0, \quad B = 0
}
zu bestimmen, von denen die erste linear, die zweite quadratisch ist.

Die Gleichungen (4) sind gleichbedeutend mit
\weq{5}{
\sigma_1 = 0, \quad \sigma_2 = 0
}
%CONT>>
